\documentclass[10pt,letterpaper]{extarticle}

% ---- Packages ----
\usepackage[margin=0.35in]{geometry}
\usepackage{multicol}
\usepackage{listings}
\usepackage{xcolor}
\usepackage{titlesec}
\usepackage{enumitem}
\usepackage{amsmath}
\usepackage[T1]{fontenc}

% ---- Color palette (ink-friendly, scannable) ----
\definecolor{kw}{HTML}{0033B3}      % keywords: deep blue
\definecolor{tp}{HTML}{1750A8}      % types: slightly lighter blue
\definecolor{cm}{HTML}{8C8C8C}      % comments: medium gray
\definecolor{str}{HTML}{067D17}     % strings: green
\definecolor{num}{HTML}{1750EB}     % numbers: blue
\definecolor{rulec}{HTML}{C8C8C8}   % section rules: light gray

% ---- Layout tweaks for density ----
\setlength{\columnsep}{0.2in}
\setlength{\columnseprule}{0.3pt}
\setlength{\parindent}{0pt}
\setlength{\parskip}{1pt}
\pagestyle{empty}

% ---- Section styling: tight, bold, ruled ----
\titleformat{\section}
  {\normalfont\bfseries\large}{}{0pt}{}[{\color{rulec}\titlerule[0.5pt]}]
\titlespacing*{\section}{0pt}{4pt}{2pt}

\titleformat{\subsection}
  {\normalfont\bfseries\small}{}{0pt}{}
\titlespacing*{\subsection}{0pt}{3pt}{0pt}

% ---- Listings setup for C++ ----
\lstdefinestyle{cpp}{
  language=C++,
  basicstyle=\ttfamily\scriptsize,
  keywordstyle=\color{kw}\bfseries,
  keywordstyle=[2]\color{tp}\bfseries,
  commentstyle=\color{cm}\itshape,
  stringstyle=\color{str},
  numberstyle=\color{num},
  showstringspaces=false,
  breaklines=true,
  breakatwhitespace=false,
  columns=fullflexible,
  keepspaces=true,
  tabsize=2,
  aboveskip=2pt,
  belowskip=1pt,
  xleftmargin=0pt,
  frame=none,
  lineskip=-0.5pt,
  morekeywords=[2]{ll,size_t,pair,pairll,vector,map,unordered_map,set,unordered_set,multiset,queue,priority_queue,stack,tuple,string,auto},
  literate={~}{{$\sim$}}1
}
\lstset{style=cpp}

% ---- Minipage-style "box" that won't break across columns ----
\newenvironment{block}[1]
  {\subsection*{#1}\nopagebreak}
  {\par\addvspace{2pt}}

\begin{document}
\begin{multicols*}{2}
\raggedright
\raggedcolumns

% ==============================================================
% HEADER
% ==============================================================
\begin{center}
{\Large\bfseries CS 392 Final Cheat Sheet}\\[-1pt]
{\scriptsize Travis Falk \textbullet{} C++ competitive programming reference}
\end{center}

\begin{lstlisting}
#include <bits/stdc++.h>
using namespace std;
using ll = long long;
using pairll = pair<ll,ll>;

int main() {
  ios_base::sync_with_stdio(false);
  cin.tie(NULL);
  // ... solve ...
  return 0;
}
\end{lstlisting}

% ==============================================================
% PART 1: PRIMITIVES / TEMPLATES
% ==============================================================
\section*{PART I \textemdash{} Primitives}

% ---- Vectors ----
\begin{block}{Vectors}
\begin{lstlisting}
// Dynamic array. [] and .at() are O(1).
// push/pop_back O(1). insert/erase mid O(N).
vector<ll> v;                  // empty
vector<ll> a(5);               // 5 zeros
vector<ll> b(5, 10);           // {10,10,10,10,10}
vector<ll> c = {5,2,8,1,9};
vector<ll> d(c.begin(), c.begin()+3);    // {5,2,8}
vector<vector<ll>> grid(n, vector<ll>(m, 0));

v.push_back(7);   v.emplace_back(7);
v.pop_back();
v.insert(v.begin()+2, 99);     // O(N)
v.erase(v.begin()+2);          // remove idx 2
v.erase(v.begin(), v.begin()+2);
v.clear(); v.resize(10); v.resize(10, -1);
fill(v.begin(), v.end(), 0);

ll x = v[0];    ll y = v.at(0);   // .at() throws
ll n = v.size();  bool e = v.empty();
\end{lstlisting}
\end{block}

% ---- Sorting & Algorithms ----
\begin{block}{Sorting \& <algorithm>}
\begin{lstlisting}
sort(v.begin(), v.end());                // asc
sort(v.begin(), v.end(), greater<ll>()); // desc
sort(v.begin(), v.end(),
     [](ll a, ll b){ return a > b; });   // custom
reverse(v.begin(), v.end());

auto it  = find(v.begin(), v.end(), 5);
ll   cnt = count(v.begin(), v.end(), 5);
ll   mn  = *min_element(v.begin(), v.end());
ll   mx  = *max_element(v.begin(), v.end());
ll   sum = accumulate(v.begin(), v.end(), 0LL);
iota(v.begin(), v.end(), 0);    // 0,1,2,3,...

// Dedupe: sort + unique + erase
sort(v.begin(), v.end());
v.erase(unique(v.begin(), v.end()), v.end());

// All permutations (must start sorted)
sort(v.begin(), v.end());
do { /* process v */ }
while (next_permutation(v.begin(), v.end()));
\end{lstlisting}
\end{block}

% ---- Binary Search / Bounds ----
\begin{block}{Binary Search \& Bounds (sorted vector)}
\begin{lstlisting}
// 4 shapes -- all O(log N):
//   first >= x :  lower_bound(x)
//   first >  x :  upper_bound(x)
//   last  <  x :  lower_bound(x), --it (skip if begin)
//   last  <= x :  upper_bound(x), --it (skip if begin)

bool found = binary_search(v.begin(), v.end(), 5);
auto lb = lower_bound(v.begin(), v.end(), 4);
auto ub = upper_bound(v.begin(), v.end(), 4);

// Once you have a bound iterator:
if (lb == v.end()) { /* no match */ }
ll val = *lb;
ll idx = lb - v.begin();   // O(1) for vector
*lb = 99;                  // mutate (may break sort!)
v.erase(lb);               // O(N), invalidates lb
\end{lstlisting}
\end{block}

% ---- Queues & PQs ----
\begin{block}{Queue, Stack, Priority Queue}
\begin{lstlisting}
queue<ll> q;  q.push(1); q.push(2);
ll f = q.front();  ll b = q.back();
q.pop();  q.size();  q.empty();

stack<ll> s;  s.push(1); s.push(2);
if (!s.empty()) { ll t = s.top(); s.pop(); }

// Max-heap (default)
priority_queue<ll> pq;
pq.push(10); pq.push(20);
ll top = pq.top(); pq.pop();          // 20

// Min-heap (3 template args required)
priority_queue<ll, vector<ll>, greater<ll>> minpq;

// PQ of pairs -- lex compare by first, then second.
priority_queue<pairll> pqp;          // max by first
priority_queue<pairll, vector<pairll>,
               greater<pairll>> pqmin;  // Dijkstra shape

// PQ with custom comparator (struct, not lambda)
struct Cmp {
  bool operator()(const pairll& a, const pairll& b)
    const { return a.second > b.second; } // min by 2nd
};
priority_queue<pairll, vector<pairll>, Cmp> pqc;
\end{lstlisting}
\end{block}

% ---- Maps ----
\begin{block}{Maps (ordered \& unordered)}
\begin{lstlisting}
// map: red-black tree, O(log N), needs <.
// unordered_map: hash, O(1) avg / O(N) worst.
map<string,ll> m;
unordered_map<string,ll> um;

m["apple"] = 5;              // creates if absent
ll v = m["apple"];           // WARNING: creates!
m.insert({"banana", 3});     // no-op if present
m.emplace("cherry", 7);
m.erase("apple");

// Safe existence (does NOT insert)
if (m.count("apple")) { }
if (m.find("apple") != m.end()) { }

for (auto& [k,v] : m) { }   // sorted for map

// Ordered-map binary search (NOT on unordered):
auto it = m.lower_bound(key);  // first >= key
auto jt = m.upper_bound(key);  // first >  key
// Once you have it:  it->first (key), it->second (val)
// it->second = x; is OK; key is const.
// next(it), prev(it) for neighbors.
// Min key: m.begin()->first; max: m.rbegin()->first.
\end{lstlisting}
\end{block}

% ---- Sets ----
\begin{block}{Sets (set, multiset, unordered\_set)}
\begin{lstlisting}
set<ll> s = {3,1,4,1,5};    // {1,3,4,5}
multiset<ll> ms = {1,1,2};  // {1,1,2} dupes OK
unordered_set<ll> us;

s.insert(2); s.erase(3);     // O(log N)
ms.erase(1);                 // removes ALL 1s
ms.erase(ms.find(1));        // removes ONE 1

bool has = s.count(2);       // 0 or 1 for set
ll freq = ms.count(1);       // real count

// Ordered binary search (NOT on unordered):
auto it = s.lower_bound(4);  // first >= 4
auto jt = s.upper_bound(4);  // first >  4
// Predecessor (last < x):
//   auto p = s.lower_bound(x);
//   if (p != s.begin()) --p;

*s.begin();    // min
*s.rbegin();   // max
// Values read-only thru iterator (would break order).
\end{lstlisting}
\end{block}

% ---- Strings ----
\begin{block}{Strings}
\begin{lstlisting}
string s = "hello";
string t(5, 'a');              // "aaaaa"
// getline(cin, s);            // whole line
s += " world";  s.push_back('!');  s.pop_back();
s.insert(5, "XYZ");
s.erase(0, 3);                 // erase 3 chars @ idx 0
s.replace(0, 5, "HELLO");
reverse(s.begin(), s.end());

string sub = s.substr(0, 4);   // (start, length)
string suf = s.substr(4);      // to end

size_t p = s.find("lo");       // or npos
size_t r = s.rfind("lo");      // last occurrence
if (p != string::npos) { }

ll   i = stoll("12345678901");
double d = stod("3.14");
string sn = to_string(42);
// <cctype>: isdigit isalpha isalnum isspace
//           islower isupper tolower toupper
\end{lstlisting}
\end{block}

% ---- Pairs & Tuples ----
\begin{block}{Pairs \& Tuples}
\begin{lstlisting}
pair<ll,string> p = {1, "hi"};
auto [a,b] = p;                // structured binding

tuple<ll,double,string> t = {1, 3.14, "Pi"};
ll i = get<0>(t);
auto [A,B,C] = t;              // preferred
tie(i, /*...*/) = t;           // unpack into vars

// Both compare LEXICOGRAPHICALLY.
vector<pairll> vp = {{3,1},{1,5},{2,2}};
sort(vp.begin(), vp.end());    // {{1,5},{2,2},{3,1}}
\end{lstlisting}
\end{block}

% ---- DSU ----
\begin{block}{Disjoint Set Union (Union-Find)}
\begin{lstlisting}
// find + unite: O(alpha(n)) ~ O(1) amortized.
struct DSU {
  vector<ll> parent, sz;
  DSU(ll n) : parent(n), sz(n, 1) {
    iota(parent.begin(), parent.end(), 0);
  }
  ll find(ll i) {
    if (parent[i] == i) return i;
    return parent[i] = find(parent[i]);   // path compr.
  }
  bool unite(ll i, ll j) {
    ll ri = find(i), rj = find(j);
    if (ri == rj) return false;
    if (sz[ri] < sz[rj]) swap(ri, rj);    // union by sz
    parent[rj] = ri;
    sz[ri] += sz[rj];
    return true;
  }
};
\end{lstlisting}
\end{block}

% ---- Segment Tree (point update, range sum) ----
\begin{block}{Segment Tree (point update, range sum)}
\begin{lstlisting}
// build O(N), update O(log N), query [l,r) O(log N).
// Iterative; 1-indexed internally; tree size = 2n.
class SegTree {
public:
  vector<ll> tree;
  ll n, sz;
  SegTree(vector<ll>& ar) {
    sz = ar.size();
    n = 1;
    while (n < sz) n <<= 1;
    tree.assign(2*n, 0);
    for (ll i = 0; i < sz; i++) tree[i+n] = ar[i];
    for (ll i = n-1; i > 0; i--)
      tree[i] = tree[2*i] + tree[2*i+1];
  }
  void update(ll i, ll val) {
    i += n;  tree[i] = val;
    while (i > 1) { i /= 2; tree[i] = tree[2*i] + tree[2*i+1]; }
  }
  ll query(ll l, ll r) {   // sum over [l, r)
    ll ans = 0;
    l += n; r += n;
    while (l < r) {
      if (l & 1) ans += tree[l++];
      if (r & 1) ans += tree[--r];
      l >>= 1;  r >>= 1;
    }
    return ans;
  }
};
\end{lstlisting}
\end{block}

% ---- Sieve ----
\begin{block}{Sieve of Eratosthenes (smallest prime factor)}
\begin{lstlisting}
// sieve[x] = 0 if x prime, else smallest prime factor.
// O(n log log n).
vector<ll> sieve(n+1, 0);
for (ll x = 2; x <= n; x++) {
  if (sieve[x]) continue;              // x is prime
  for (ll u = 2*x; u <= n; u += x)
    if (!sieve[u]) sieve[u] = x;
}
// To factor k: while (sieve[k]) { p = sieve[k]; k/=p; }
//              then k itself is a prime factor (if >1).
\end{lstlisting}
\end{block}

% ---- BFS ----
\begin{block}{BFS (unweighted shortest path)}
\begin{lstlisting}
// O(V + E). dist[i] = hops from src, -1 if unreachable.
vector<ll> bfs(ll src, ll n, vector<vector<ll>>& adj) {
  vector<ll> dist(n+1, -1);
  queue<ll> q;
  dist[src] = 0;  q.push(src);
  while (!q.empty()) {
    ll u = q.front(); q.pop();
    for (ll v : adj[u]) if (dist[v] == -1) {
      dist[v] = dist[u] + 1;
      q.push(v);
    }
  }
  return dist;
}
\end{lstlisting}
\end{block}

% ---- DFS ----
\begin{block}{DFS (iterative \& recursive)}
\begin{lstlisting}
// Iterative: component detection, reachability.
vector<bool> dfs_it(ll src, ll n, vector<vector<ll>>& adj) {
  vector<bool> vis(n+1, false);
  stack<ll> st;  st.push(src);
  while (!st.empty()) {
    ll u = st.top(); st.pop();
    if (vis[u]) continue;
    vis[u] = true;
    for (ll v : adj[u]) if (!vis[v]) st.push(v);
  }
  return vis;
}

// Recursive: trees, backtracking, coloring.
vector<bool> vis;   // global to avoid passing
void dfs(ll u, vector<vector<ll>>& adj) {
  vis[u] = true;
  for (ll v : adj[u]) if (!vis[v]) dfs(v, adj);
}
// Usage:  vis.assign(n+1, false);  dfs(src, adj);
\end{lstlisting}
\end{block}

% ---- Dijkstra ----
\begin{block}{Dijkstra (non-negative weights)}
\begin{lstlisting}
// O((V+E) log V). adj[u] = {(v, w), ...}.
vector<ll> dijkstra(ll src, ll n,
                    vector<vector<pairll>>& adj) {
  const ll INF = 1e18;
  vector<ll> dist(n+1, INF);
  priority_queue<pairll, vector<pairll>,
                 greater<pairll>> pq;     // min-heap
  dist[src] = 0;
  pq.push({0, src});
  while (!pq.empty()) {
    auto [d, u] = pq.top(); pq.pop();
    if (d > dist[u]) continue;            // stale
    for (auto [v, w] : adj[u]) {
      if (dist[u] + w < dist[v]) {
        dist[v] = dist[u] + w;
        pq.push({dist[v], v});
      }
    }
  }
  return dist;
}
\end{lstlisting}
\end{block}

% ---- DP Templates ----
\begin{block}{DP: 0/1 Knapsack (1D)}
\begin{lstlisting}
// dp[j] = max value at capacity j.
// Items OUTER, capacity BACKWARD (prevents re-use).
vector<ll> dp(W+1, 0);
for (ll i = 0; i < n; i++) {
  for (ll j = W; j >= weight[i]; j--) {
    dp[j] = max(dp[j], dp[j-weight[i]] + value[i]);
  }
}
cout << dp[W];
\end{lstlisting}
\end{block}

\begin{block}{DP: Unbounded Knapsack (1D)}
\begin{lstlisting}
// Each item takeable unlimited times.
// Items OUTER, capacity FORWARD (allows re-use).
vector<ll> dp(W+1, 0);
for (ll i = 0; i < n; i++) {
  for (ll j = weight[i]; j <= W; j++) {
    dp[j] = max(dp[j], dp[j-weight[i]] + value[i]);
  }
}
\end{lstlisting}
\end{block}

\begin{block}{DP: Bounded Knapsack (binary decomp \textrightarrow{} 0/1)}
\begin{lstlisting}
// copies[i] copies of item i. Decompose into bundles
// of sizes 1,2,4,...,rem which sum to copies[i].
// Any count 0..copies[i] = subset sum of bundles.
vector<pairll> bundles;  // (weight, value)
for (ll i = 0; i < n; i++) {
  ll k = copies[i];
  for (ll t = 1; t <= k; t *= 2) {
    bundles.push_back({t*weight[i], t*value[i]});
    k -= t;
  }
  if (k > 0)
    bundles.push_back({k*weight[i], k*value[i]});
}
// Then standard 0/1 knapsack over bundles.
vector<ll> dp(W+1, 0);
for (auto [w, v] : bundles) {
  for (ll j = W; j >= w; j--) {
    dp[j] = max(dp[j], dp[j-w] + v);
  }
}
\end{lstlisting}
\end{block}

\begin{block}{DP: Counting (ordered composition)}
\begin{lstlisting}
// # ways to make sum n using parts from S, order matters.
// e.g. dice: S = {1,2,3,4,5,6}.  1+2 != 2+1.
// Target OUTER, parts INNER.
const ll MOD = 1000000007;
vector<ll> dp(n+1, 0);
dp[0] = 1;
for (ll s = 1; s <= n; s++) {
  for (ll k : {1,2,3,4,5,6}) {
    if (s - k >= 0)
      dp[s] = (dp[s] + dp[s-k]) % MOD;
  }
}
cout << dp[n];
\end{lstlisting}
\end{block}

\begin{block}{DP: Counting (unordered -- coin change)}
\begin{lstlisting}
// # ways to make sum n using coins in C, order irrelevant.
// 1+2 and 2+1 counted ONCE.
// Coins OUTER, target INNER.  (swap vs ordered version)
vector<ll> dp(n+1, 0);
dp[0] = 1;
for (ll c : coins) {
  for (ll s = c; s <= n; s++) {
    dp[s] += dp[s - c];
  }
}
\end{lstlisting}
\end{block}

\begin{block}{DP: Longest Common Subsequence}
\begin{lstlisting}
// dp[i][j] = LCS length of a[0..i) and b[0..j).
ll n = a.size(), m = b.size();
vector<vector<ll>> dp(n+1, vector<ll>(m+1, 0));
for (ll i = 1; i <= n; i++) {
  for (ll j = 1; j <= m; j++) {
    if (a[i-1] == b[j-1])
      dp[i][j] = dp[i-1][j-1] + 1;
    else
      dp[i][j] = max(dp[i-1][j], dp[i][j-1]);
  }
}
cout << dp[n][m];
\end{lstlisting}
\end{block}

\begin{block}{DP: Edit Distance (Levenshtein)}
\begin{lstlisting}
// dp[i][j] = min edits to make a[0..i) == b[0..j).
vector<vector<ll>> dp(n+1, vector<ll>(m+1));
for (ll i = 0; i <= n; i++) dp[i][0] = i;  // deletions
for (ll j = 0; j <= m; j++) dp[0][j] = j;  // insertions
for (ll i = 1; i <= n; i++) {
  for (ll j = 1; j <= m; j++) {
    if (a[i-1] == b[j-1])
      dp[i][j] = dp[i-1][j-1];
    else
      dp[i][j] = 1 + min({dp[i-1][j],       // delete
                          dp[i][j-1],       // insert
                          dp[i-1][j-1]});   // replace
  }
}
\end{lstlisting}
\end{block}

% ---- DP direction cheat sheet ----
\begin{block}{DP Knapsack-Family Direction Table}
\begin{lstlisting}
// Problem         | Outer  | Inner     | Note
// --------------- | ------ | --------- | -----------------
// 0/1 knapsack    | items  | cap BACK  | each used <= 1x
// Unbounded knap  | items  | cap FWD   | unlimited uses
// Bounded knap    | bundles| cap BACK  | decompose first
// Count ordered   | target | parts     | 1+2 != 2+1
// Count unordered | coins  | target    | 1+2 == 2+1
//
// Mental model when collapsing 2D -> 1D:
//   "previous row" lives in not-yet-overwritten cells.
//   BACKWARD iter: reads old row's smaller index.
//   FORWARD iter: reads THIS row's smaller index.
\end{lstlisting}
\end{block}


\end{multicols*}
\end{document}
